% !TeX root = ../thuthesis-example.tex

\chapter{总结与展望}

\section{工作总结}

本文围绕生成式推荐中语义ID的构建问题，提出了COMET（Collaborative-guided Multimodal Encoding Tokenizer）方法。主要工作和贡献总结如下：

（1）针对现有语义ID构建方法模态利用单一、协同过滤信号融合间接的问题，提出了CF-as-Query跨注意力融合机制。该机制以协同过滤嵌入为查询，以文本多token投影和图像嵌入构成的键值序列为键值对，通过堆叠的跨注意力和前馈网络自适应地从内容特征中提取推荐相关信息。与LETTER的辅助损失方式相比，COMET将协同过滤信号从损失层面的间接引导提升为架构层面的直接融合，无需引入额外的对比损失函数。

（2）提出了多目标重建策略，通过三个独立的解码器分支以加权方式同时重建文本、图像和协同过滤嵌入，促进多模态信息在64维隐空间中的均衡编码。配合余弦退火学习率调度和训练时Sinkhorn均匀化等训练策略，进一步提升了语义ID的质量。

（3）搭建了从数据处理、多模态嵌入提取、语义ID构建到生成式推荐的完整实验流水线，实现了RQ-KMeans、RQ-VAE、LETTER和COMET四种语义ID构建方法在TIGER和OneRec两种生成式推荐模型上的公平对比。实验结果表明，COMET在所有设置下均取得了最优的推荐性能。

\section{未来展望}

尽管COMET取得了较好的实验结果，仍有以下方向值得进一步探索：

（1）\textbf{更多模态的融入}。当前COMET融合了文本、图像和协同过滤三种模态。在实际应用中，物品还可能具有其他模态信息（如视频、评论文本情感、价格区间等），跨注意力架构天然支持扩展更多的Key-Value token，未来可以探索更丰富的模态融合。

（2）\textbf{端到端联合训练}。当前COMET的语义ID构建和下游生成模型是分阶段训练的。将两者端到端联合优化可能进一步提升性能，使语义ID的构建直接面向推荐目标进行优化。

（3）\textbf{更大规模数据集验证}。当前实验仅在Amazon Beauty数据集上进行。未来应在更多领域和更大规模的数据集上验证COMET的泛化能力。
